\chapter{Conclusion}
\label{chap:conclusion}

This project developed and evaluated a multimodal emotion recognition system that classifies six emotion categories from video input. The system processes three complementary streams — text transcripts, acoustic features, and facial visual features — and fuses them through an attention bottleneck mechanism. The work addressed the full research pipeline: literature analysis, dataset preparation, architecture design, training under class imbalance, quantitative evaluation, and integration into a functional prototype backend.

The proposed architecture combines three modality-specific encoders with a bottleneck fusion module. The text branch uses frozen BERT-base-uncased and averages the last four hidden layers to build word-level representations. This avoids overfitting that occurs when fine-tuning 110 million parameters on a dataset of 13\,934 training samples. The audio and visual branches use Conv1D projections with positional encoding to bring COVAREP and OpenFace features into a shared 128-dimensional space. The fusion module uses 16 learnable bottleneck tokens as a compact communication channel. All cross-modal interaction passes through these tokens rather than through full pairwise attention between sequence positions, which keeps the complexity at $\mathcal{O}(T \cdot N_B)$ instead of $\mathcal{O}(T^2)$. Learnable modal weights control how much each stream contributes to the bottleneck state. Separate residual connections and LayerNorm instances guard each attention step. Auxiliary classification heads on each stream and modality dropout with probability 0.05 prevent the model from relying only on the strongest modality.

The experimental results confirm that bottleneck attention outperforms both implemented baselines. Baseline 1, which combines BERT with Conv1D audio encoding and a BiLSTM visual encoder using late fusion, reached Avg.~F1 of 0.4412. Baseline 2, based on MulT-style pairwise cross-modal attention, improved this to 0.4631. The proposed bottleneck model achieved 0.4854 with the default decision threshold and 0.4942 after per-class threshold tuning on the validation set. This represents a 4.42 percentage point gain over Baseline 1 and shows that structured bottleneck communication adds value beyond both late fusion and unconstrained cross-attention.

Recognition of rare emotion categories improved substantially. Fear F1 rose from 0.1834 in Baseline 1 to 0.3122 at the default threshold and 0.3451 after tuning. Surprise F1 increased from 0.1902 to 0.3214 after tuning. These gains come from a combination of class-specific positive weights in the loss function, modality dropout regularisation, and per-class threshold selection on the validation set. Among published systems that use compact acoustic and visual descriptors, the proposed model surpasses MulT and approaches XMBT in Avg.~F1.

The ablation study identified which design choices contributed most. Frozen BERT with averaged last four layers provided a stronger and more stable text representation than online fine-tuning. Auxiliary losses ensured that audio and visual branches retained task-relevant information throughout training. The Semantic Enhancement Module did not improve results in this setting, which shows that components from other architectures do not always transfer directly when the feature space differs.

The project also delivered a functional prototype. The FastAPI backend accepts raw video, extracts audio with ffmpeg, computes COVAREP features with opensmile, detects facial action units with py-feat, transcribes speech with faster-whisper, and returns emotion probabilities through a REST API. The system is currently a research prototype. Preprocessing, especially transcription and facial feature extraction, dominates inference time and limits real-time use.

The main limitations are the use of handcrafted acoustic and visual features, evaluation on a single English-language dataset, and the absence of cross-dataset testing. These constrain both the representational capacity of the model and the generalisability of the results.

Future work should replace COVAREP with pretrained speech encoders such as wav2vec2 or HuBERT, and OpenFace descriptors with a neural facial expression model. Additional directions include deeper bottleneck layers, larger token sets, focal loss for rare classes, and cross-dataset evaluation. On the engineering side, optimised preprocessing and GPU deployment would move the system closer to interactive use.

This project confirms that attention bottleneck fusion is an effective and computationally efficient strategy for multimodal emotion recognition. The developed system improves over both baselines, handles class imbalance more effectively, and provides a solid foundation for future work with stronger modality encoders.